% TEMPLATE-MILC-v3.tex - plantilla maestra UNICA de las guias MILC v3.
%
% La usan las cuatro familias de guias, cada una con su driver:
%   · Grados 6-11          -> scripts/build-guias-g11.py
%   · Bebras               -> scripts/build-guias-bebras.py
%   · Semillero            -> scripts/build-guias-semillero.py
%   · Territorio interior  -> scripts/build-guias-territorio-interior.py
%
% Antes habia tres copias casi identicas (~400 lineas iguales, 6 distintas):
% cada mejora habia que aplicarla tres veces, y en la practica no se hacia
% (el motor de recursos vivia solo en dos de ellas). Lo que varia entre
% familias esta parametrizado en 6 placeholders que cada driver llena
% (se nombran aqui SIN su sintaxis real: el builder reemplaza por texto plano,
% tambien dentro de los comentarios, y un valor multilinea rompe el preambulo):
%
%   HEADER_IZQ        encabezado izquierdo de pagina
%   HEADER_DER        encabezado derecho de pagina
%   PORTADA_ETIQUETA  rotulo sobre el numero grande (GRADO/MOMENTO/MODULO)
%   PORTADA_SERIE     linea de serie de la portada
%   PORTADA_ID        identificador de la guia en la portada
%   PORTADA_CIRCULO   el circulito de grado (°), o vacio si no aplica
%   PDF_TITULO        titulo en texto plano (sin \textbf ni comillas TeX) para
%                     los metadatos del PDF (/Title); lo produce lib_xelatex.texto_pdf
%
% Composicion (auditoria 2026-09-03): las bandas de estacion piden espacio
% (\Needspace*) y prohiben el salto justo despues (\nopagebreak), asi no quedan
% huerfanas al pie; las cajas largas (softbox, drawbox, infoband, puentebox) son
% `breakable`, asi una caja de media pagina ya no arrastra todo a la siguiente
% dejando la anterior al 40 %. Todo texto sobre mostaza o verde va en milcNegro
% (blanco sobre #E5B400 da 1,93:1 y sobre #58A923 2,95:1; ninguno pasa WCAG AA).
% El resto de claves las documenta content/guias/_SCHEMA.md. El script valida
% que no queden placeholders sin reemplazar antes de escribir el archivo final.

% Etiquetado PDF (latex-lab): /Lang, estructura y texto alternativo de las
% figuras, para lectores de pantalla. Debe ir ANTES de \documentclass.
\DocumentMetadata{lang=es-CO,tagging=on}
\documentclass[11pt,letterpaper]{article}
% headheight=24pt: la cabecera derecha lleva dos lineas (identidad de la guia
% y estacion actual). headsep se acorta en la misma medida para que el bloque
% de texto quede exactamente donde estaba con headheight=12pt.
\usepackage[margin=2.0cm,headheight=24pt,headsep=13pt]{geometry}
\usepackage[table]{xcolor}
\usepackage{fontspec}
\usepackage[spanish]{babel}
\usepackage{tikz}
% Librerías TikZ para los diagramas de las guías (bloque `recursos:`):
%  · babel  → babel en español vuelve activos caracteres como «>», y TikZ los usa
%             en sus opciones (>=stealth, ->). Sin esta librería, cualquier
%             diagrama con flechas revienta con «\language@active@arg> has an
%             extra }». Debe ir ANTES de que se dibuje nada.
%  · positioning        → sintaxis «below left=2pt and 3pt».
%  · decorations.pathreplacing → llaves ({brace}) para señalar rangos.
\usetikzlibrary{babel,positioning,decorations.pathreplacing}
\usepackage{graphicx}
% Circuitos: el trabajo de electrónica se hace hoy con micro:bit y sus sensores
% integrados (MakeCode, sin cableado), así que no se necesitan esquemáticos.
% Si algún día entran montajes con protoboard, basta descomentar esta línea:
% los diagramas .tex/.tikz declarados en `recursos:` ya se incrustan solos.
% \usepackage{circuitikz}
\usepackage[most]{tcolorbox}
\usepackage{tabularx}
\usepackage{array}
\usepackage{fancyhdr}
\usepackage{titlesec}
\usepackage[document]{ragged2e}  % bandera a la derecha en todo el cuerpo
\usepackage{enumitem}
\usepackage{microtype}
\usepackage{etoolbox}
\usepackage{needspace}
% hyperref al final del preambulo: metadatos del PDF (titulo, autor, idioma,
% familia), marcadores por estacion (\pdfbookmark) y enlaces sin recuadro.
\usepackage[hidelinks,
  pdftitle={Fórmulas compuestas — el orden de las operaciones y los paréntesis que lo dicen},
  pdfauthor={PhD. Álvaro Cárdenas Orozco},
  pdfsubject={Tecnología e Informática · Grado 8 · Guía 6 · MILC v3},
  pdflang={es-CO},
  bookmarksopen=true,bookmarksnumbered=false]{hyperref}

% Helvetica Neue no trae la flecha U+2192, asi que cada «→» escrito literalmente
% en el YAML se compilaba a NADA, en silencio (462 flechas en 61 guias: rutas del
% tipo «paleta → tipografia → lockup» salian rotas). Mapeamos el caracter a su
% version matematica, que si tiene glifo. Asi el autor puede seguir escribiendo
% «→» comodamente en el YAML.
\usepackage{newunicodechar}
\newunicodechar{→}{\ensuremath{\rightarrow}}
% Mismo problema con los simbolos de marca y vinetas: el «✓» de «una palomita ✓»
% salia como cuadrito vacio en el PDF de todas las guias que lo usaban.
\usepackage{pifont}
\newunicodechar{✓}{\ding{51}}
\newunicodechar{✗}{\ding{55}}
\newunicodechar{✔}{\ding{52}}
\newunicodechar{•}{\textbullet}
\newunicodechar{≥}{\ensuremath{\geq}}
\newunicodechar{≤}{\ensuremath{\leq}}

\setmainfont{Helvetica Neue}
\setsansfont{Helvetica Neue}
\setlength{\parindent}{0pt}
\setlength{\parskip}{6pt}
% Sin particion de palabras: en 6.º-7.º una palabra cortada por guion al final
% de linea («Comuni-cacion») frena la lectura mucho mas que una linea corta.
\hyphenpenalty=10000
\exhyphenpenalty=10000
\pretolerance=2000
\tolerance=3000
\setlength{\emergencystretch}{3em}
\linespread{1.10}
\renewcommand{\arraystretch}{1.25}
\setlist{leftmargin=5mm,itemsep=1mm,topsep=1mm}
\newcolumntype{Y}{>{\RaggedRight\arraybackslash}X}

\definecolor{milcMagenta}{HTML}{E6007E}
\definecolor{milcVino}{HTML}{5A0038}
\definecolor{milcTurquesa}{HTML}{0093A5}
\definecolor{milcVerde}{HTML}{58A923}
\definecolor{milcMostaza}{HTML}{E5B400}
\definecolor{milcOcre}{HTML}{B86600}
\definecolor{milcNegro}{HTML}{241816}
\definecolor{milcGris}{HTML}{F7F7F4}
\definecolor{milcLinea}{HTML}{E8E2DD}
\definecolor{milcGradePrimary}{HTML}{FF6600}
\definecolor{milcGradeDark}{HTML}{9A3E00}
\definecolor{milcGradeSoft}{HTML}{FFE4D0}
\definecolor{milcGradeAccent}{HTML}{CFE7FF}
\definecolor{milcGradeText}{HTML}{FFFFFF}
\color{milcNegro}

\pagestyle{fancy}
\fancyhf{}
% Cabecera de dos lineas: identidad de la guia arriba y, a la derecha y en
% gris, la estacion en curso (\markright desde \milcestacion). Antes de la
% primera estacion la segunda linea va vacia; el \strut mantiene la altura
% para que la linea de arriba no baile entre paginas.
\lhead{\small Tecnología e Informática · Grado 8\\[-1pt]{\footnotesize\strut}}
\rhead{\small Guía 6 · MILC v3\\[-1pt]{\footnotesize\strut\color{milcNegro!65}\rightmark}}
\cfoot{\small\thepage}
\renewcommand{\headrulewidth}{0pt}

% Sobre mostaza (#E5B400) y verde (#58A923) el texto va en milcNegro; sobre el
% resto de la paleta, en blanco. Se usa dentro de fonttitle/fontupper, donde un
% \color posterior manda sobre coltitle.
\newcommand{\milctintasobre}[1]{%
  \ifboolexpr{test{\ifstrequal{#1}{milcMostaza}} or test{\ifstrequal{#1}{milcVerde}}}%
    {\color{milcNegro}}{\color{white}}}

\newtcolorbox{titlebox}[1]{
  enhanced,colback=#1,colframe=#1,arc=7mm,boxrule=0pt,
  left=7mm,right=7mm,top=3mm,bottom=3mm,
  before skip=8pt,after skip=8pt,
  fontupper=\sffamily\bfseries\Large\color{white}
}

\newtcolorbox{softbox}[3]{
  enhanced,breakable,pad at break=3mm,
  colback=#2,colframe=#1,borderline west={4pt}{0pt}{#1},
  arc=2mm,boxrule=.4pt,left=4mm,right=4mm,top=3mm,bottom=3mm,
  before skip=6pt,after skip=8pt,title={#3},coltitle=white,
  colbacktitle=#1,fonttitle=\sffamily\bfseries\milctintasobre{#1},fontupper=\RaggedRight,
  attach boxed title to top left={xshift=3mm,yshift=-2mm},
  boxed title style={arc=2mm, boxrule=0pt}
}

\newtcolorbox{drawbox}[2]{
  enhanced,breakable,pad at break=4mm,
  colback=white,colframe=#1,arc=4mm,boxrule=1pt,
  left=5mm,right=5mm,top=4mm,bottom=4mm,before skip=8pt,after skip=8pt,
  title={#2},coltitle=white,colbacktitle=#1,fonttitle=\sffamily\bfseries\milctintasobre{#1},
  fontupper=\RaggedRight,
  attach boxed title to top left={xshift=4mm,yshift=-2mm},
  boxed title style={arc=3mm, boxrule=0pt}
}

\newtcolorbox{infoband}[1]{
  enhanced,breakable,pad at break=2mm,colback=#1!8,colframe=#1!50,arc=2mm,boxrule=.5pt,
  left=4mm,right=4mm,top=2mm,bottom=2mm,before skip=4pt,after skip=4pt,
  fontupper=\small\RaggedRight
}

% Puente narrativo entre secciones (contrato editorial Conéctate).
% Da continuidad: "Acabas de X. Ahora vas a Y porque Z."
% Visualmente: banda izquierda en color del grado, texto en itálica pequeña.
\newtcolorbox{puentebox}{
  enhanced,breakable,pad at break=2mm,colback=milcGradeSoft,colframe=milcGradeDark,
  borderline west={3pt}{0pt}{milcGradeDark},
  arc=0mm,boxrule=0pt,left=5mm,right=4mm,top=2.5mm,bottom=2.5mm,
  before skip=4pt,after skip=8pt,
  fontupper=\itshape\small\color{milcNegro}\RaggedRight
}
% La flecha va en modo matematico a proposito: Helvetica Neue NO tiene el glifo
% U+2192, asi que \textrightarrow se compilaba a NADA («Missing character: There
% is no → in font Helvetica Neue») y el puente salia sin flecha. En modo
% matematico la flecha viene de la fuente matematica y si se dibuja.
\newcommand{\puente}[1]{%
  \begin{puentebox}%
    \textcolor{milcGradeDark}{$\rightarrow$}\hspace{2mm}#1%
  \end{puentebox}%
}

\newcommand{\linead}[1]{\noindent\color{milcLinea}\rule{#1}{0.5pt}\color{milcNegro}}
\newcommand{\respuesta}[1]{\par\vspace{2mm}\foreach \n in {1,...,#1}{\noindent\color{milcLinea}\rule{\linewidth}{0.5pt}\par\vspace{5mm}}\color{milcNegro}}
\newcommand{\checkbox}{\textcolor{milcGradeDark}{\fbox{\phantom{x}}}\hspace{5pt}}

% ─── Listas aireadas para las guías socioemocionales ────────────────────
% Componer las verificaciones como «\checkbox texto\\» deja el texto que salta
% de línea debajo del cuadrito y apelmaza el bloque. Estos entornos alinean la
% continuación y airean los ítems. Los usa Territorio Interior; cualquier
% familia puede hacerlo.
\newlist{checklista}{itemize}{1}
\setlist[checklista]{
  label=\textcolor{milcGradeDark}{\fbox{\phantom{x}}},
  leftmargin=9mm, labelsep=3.2mm, itemsep=2.4mm,
  topsep=2.5mm, parsep=0pt, partopsep=0pt, before=\RaggedRight
}
\newlist{pasoslista}{enumerate}{1}
\setlist[pasoslista]{
  label=\textcolor{milcGradeDark}{\sffamily\bfseries\arabic*.},
  leftmargin=8mm, labelsep=3mm, itemsep=2.8mm,
  topsep=1mm, parsep=0pt, partopsep=0pt, before=\RaggedRight
}

\newcommand{\phasecard}[4]{
\begin{tcolorbox}[enhanced,colback=#3,colframe=#2,arc=3mm,boxrule=.5pt,height=3.05cm,valign=center,halign=center,left=1.5mm,right=1.5mm,top=2mm,bottom=2mm]
{\sffamily\bfseries\fontsize{22}{24}\selectfont\textcolor{#2}{#1}}\par
{\sffamily\bfseries\small #4}
\end{tcolorbox}
}

% ── Piezas del rediseno 2026 (legibilidad 6.º-11.º) ───────────────────
% Banda de estacion: reemplaza el titlebox macizo. Titulo a la izquierda,
% dato practico (tiempo/modalidad) a la derecha, en una sola linea.
% Nunca huerfana: pide 9 lineas libres (o salta de pagina rellenando con
% \vfill) y prohibe el salto justo despues de la banda. Nueve y no seis:
% la caja partible que sigue necesita ~80pt (titulo adosado, relleno y dos
% lineas) para arrancar en la misma pagina; con menos, tcolorbox la manda
% entera a la siguiente y la banda se queda sola. El penalty 10000 va
% en `after`, ANTES del espacio posterior: un `after skip` (aunque sea 0pt)
% es un \vskip tras la caja, y ese glue es un punto de corte legal que un
% \nopagebreak escrito despues de \end{tcolorbox} ya no cierra. Cada banda
% deja marcador PDF y actualiza la cabecera derecha.
\newcounter{milcestacion}
\newcommand{\milcestacion}[2]{%
\Needspace*{9\baselineskip}%
\stepcounter{milcestacion}%
\pdfbookmark[1]{#2}{milcestacion\themilcestacion}%
\markright{#2}%
\begin{tcolorbox}[enhanced,colback=#1,colframe=#1,arc=2mm,boxrule=0pt,
  left=5mm,right=5mm,top=2.6mm,bottom=2.6mm,before skip=11pt,
  after={\par\nopagebreak\vskip7pt}]
{\sffamily\bfseries\fontsize{15}{18}\selectfont\milctintasobre{#1}#2}
\end{tcolorbox}}

% Tarjeta de paso numerada: sustituye las tablas «Pilar / Aplicacion», que
% eran formato de consulta docente, no de estudiante. El numero va en un
% circulo sobre el borde para que la secuencia se lea de un vistazo.
\newcommand{\milcpaso}[3]{%
\Needspace{4\baselineskip}%
\begin{tcolorbox}[enhanced,colback=white,colframe=milcLinea,arc=1.5mm,boxrule=.9pt,
  left=14mm,right=4mm,top=3mm,bottom=3mm,before skip=4pt,after skip=4pt,
  fontupper=\RaggedRight,
  overlay={\node[anchor=north west,circle,fill=#2,text=white,inner sep=0pt,
    minimum size=7.4mm,font=\sffamily\bfseries\small]
    at ([xshift=3.6mm,yshift=-3.2mm]frame.north west) {#1};}]
#3
\end{tcolorbox}}

% Casilla marcable de tamano util con lapiz (la anterior era \fbox{\phantom{x}},
% de alto variable segun el texto que la seguia).
\newcommand{\milcmarca}{\framebox[3.8mm]{\rule{0pt}{3.4mm}}}

% Fila marcable de lista suelta (checks de la estacion 1).
\newcommand{\milccheckrow}[1]{%
\par\vspace{1.8mm}\noindent
\makebox[6.4mm][l]{\raisebox{-.55mm}{\textcolor{milcNegro}{\milcmarca}}}%
\parbox[t]{\dimexpr\linewidth-6.4mm\relax}{\RaggedRight #1}\par}

% Celda marcable dentro de la rubrica (tabla de autoevaluacion). Misma
% sangria francesa, pero pensada para vivir dentro de una columna X.
\newcommand{\milcopcion}[1]{%
\makebox[5.2mm][l]{\raisebox{-.55mm}{\textcolor{milcNegro}{\milcmarca}}}%
\parbox[t]{\dimexpr\linewidth-5.2mm\relax}{\RaggedRight #1}}

% ── Recursos visuales de la guía ──────────────────────────────────────
% Declarados en el bloque `recursos:` del YAML y emitidos por el builder.
% \guiaFigura   → imagen rasterizada o PDF (foto, carta celeste, montaje).
% \guiaDiagrama → fuente TikZ/circuitikz (.tex) incrustada como vector:
%                 es la vía para circuitos y esquemáticos editables.
% [#1] = texto alternativo (alt) · #2 = ruta relativa al .tex · #3 = pie
% El alt llega al PDF etiquetado (/Alt) y lo lee el lector de pantalla.
\newcommand{\guiaFigura}[3][]{%
\begin{center}
\includegraphics[width=0.88\linewidth,height=0.38\textheight,keepaspectratio,alt={#1}]{#2}\\[1.5mm]
{\footnotesize\itshape\textcolor{milcNegro!75}{#3}}
\end{center}
}

\newcommand{\guiaDiagrama}[2]{%
\begin{center}
\input{#1}\\[1.5mm]
{\footnotesize\itshape\textcolor{milcNegro!75}{#2}}
\end{center}
}

\begin{document}
\thispagestyle{empty}

% ── Portada ───────────────────────────────────────────────────────────
% Antes: un rectangulo de color a pagina completa. Se veia bien en pantalla,
% pero en la fotocopiadora del colegio se comia el toner y salia gris sucio.
% Ahora la banda ocupa el tercio superior y el resto va en blanco.
\begin{tikzpicture}[remember picture,overlay]
  \path[fill=milcGradePrimary]
    (current page.north west) rectangle ([yshift=-10.6cm]current page.north east);

  \node[anchor=north west,text=milcGradeText] at ([xshift=2.0cm,yshift=-1.55cm]current page.north west)
    {\sffamily\bfseries\fontsize{12}{14}\selectfont GRADO};
  \node[anchor=north west,text=milcGradeText,opacity=.85] at ([xshift=2.0cm,yshift=-2.35cm]current page.north west)
    {\sffamily\bfseries\fontsize{8.5}{10}\selectfont SERIE GUÍAS · TECNOLOGÍA E INFORMÁTICA};
  \node[anchor=north east,text=milcGradeText,opacity=.85,align=right] at ([xshift=-2.0cm,yshift=-1.55cm]current page.north east)
    {\sffamily\bfseries\fontsize{9}{12}\selectfont I.E. Sor María Juliana\\Cartago, Valle del Cauca};

  \node[anchor=west,text=milcGradeText] (gradonum) at ([xshift=1.85cm,yshift=-7.1cm]current page.north west)
    {\sffamily\bfseries\fontsize{112}{112}\selectfont 8};
  % El circulito de grado (°) solo aplica a las guias de grado («11°»); en
  % Bebras y Semillero el numero grande es un momento/modulo, no un grado.
  % Se posiciona relativo al nodo `gradonum`, no en coordenadas absolutas.
  \draw[milcGradeText,line width=1.6pt]
    ([xshift=2.0mm,yshift=-4.5mm]gradonum.north east) circle (.15cm);

  \node[anchor=west,text=milcGradeText,text width=11.4cm,align=left] at ([xshift=1.5cm,yshift=0.5cm]gradonum.east)
    {
     {\sffamily\bfseries\fontsize{9}{11}\selectfont\MakeUppercase{Período 1 · Análisis de datos con phronesis}}\\[3mm]
     {\sffamily\bfseries\fontsize{21}{25}\selectfont\setlength{\baselineskip}{27pt}El orden de\\[4pt]las operaciones}\\[3.5mm]
     {\sffamily\bfseries\fontsize{15}{18}\selectfont Guía 6-8-TIC}
    };
\end{tikzpicture}

\vspace*{9.4cm}
\vfill

{\sffamily\bfseries\fontsize{8}{10}\selectfont\textcolor{milcGradeDark}{LO QUE TE LLEVAS HOY}}

\begin{tcolorbox}[enhanced,colback=white,colframe=milcNegro,arc=2mm,boxrule=1.6pt,
  left=5mm,right=5mm,top=4mm,bottom=4mm,before skip=3pt,after skip=12pt,
  fontupper=\RaggedRight\fontsize{11}{15.5}\selectfont]
Una hoja de Excel con cinco fórmulas compuestas reales (nota final ponderada, precio con IVA y descuento, promedio sin el menor valor, entre otras), escritas con paréntesis que dejan claro el orden. Además, tres fórmulas mal escritas a propósito, corregidas y explicadas en una línea, y una prueba final con un cálculo de tu vida real.
\end{tcolorbox}

\begin{tcolorbox}[enhanced,colback=milcGris,colframe=milcGris,arc=2mm,boxrule=0pt,
  left=5mm,right=5mm,top=3.5mm,bottom=3.5mm,after skip=12pt]
{\sffamily\bfseries\fontsize{8}{10}\selectfont\textcolor{milcNegro!55}{ESTA GUÍA ES DE}}\\[3mm]
\begin{tabularx}{\linewidth}{X p{3.2cm} p{3.2cm}}
\linead{\linewidth} & \linead{3.0cm} & \linead{3.0cm}\\[-1mm]
{\fontsize{8.5}{10}\selectfont\textcolor{milcNegro!55}{Nombre completo}} &
{\fontsize{8.5}{10}\selectfont\textcolor{milcNegro!55}{Grupo}} &
{\fontsize{8.5}{10}\selectfont\textcolor{milcNegro!55}{Fecha}}\\
\end{tabularx}
\end{tcolorbox}

\vfill

\noindent\textcolor{milcLinea}{\rule{\linewidth}{1pt}}\\[2mm]
{\sffamily\bfseries\fontsize{9.5}{12}\selectfont Autor: Dr. Álvaro Cárdenas Orozco}\\[1mm]
{\sffamily\fontsize{8.5}{11}\selectfont\textcolor{milcNegro!55}{Metodología MILC · apertura ancestral · pensamiento computacional · triángulo Dussel–estoicismo–Floridi}}

\newpage

\section*{Apertura: Fórmulas compuestas --- el orden de las operaciones y los paréntesis que lo dicen}

\begin{softbox}{milcGradeDark}{milcGradeSoft}{Saber ancestral}
En las selvas del Chocó, en el litoral del río San Juan, las mujeres wounaan tejen canastos con la fibra de la palma de werregue. La técnica es el \textbf{rollo en espiral}: se empieza por el centro y cada vuelta se cose sobre la anterior, hasta que el canasto cierra. Los diseños cuentan la vida cotidiana, la naturaleza y las creencias del pueblo (Universidad del Rosario, 2023). Fíjate en el orden. No se puede tejer la vuelta de afuera antes que la del centro. No se puede volver a la del medio cuando ya cerraste la de arriba. Cada vuelta depende de la anterior, y el diseño solo aparece si el orden se respeta. Eso es una fórmula compuesta: varias operaciones que solo dan el resultado correcto en un orden. La \textbf{cara de exclusión}: la técnica se enseña hoy en talleres universitarios de Bogotá. La página del taller nombra al pueblo y a la técnica, pero a ninguna tejedora. Hoy vas a escribir fórmulas donde el orden importa tanto como en el canasto, y vas a usar los paréntesis para dejarlo claro.\par\smallskip{\footnotesize\textcolor{milcNegro!70}{\textbf{Fuente.} Universidad del Rosario, Facultad de Creación. (2023). \emph{Saberes artesanales: taller de tejeduría en werregue}.}}
\end{softbox}

\begin{softbox}{milcTurquesa}{milcGris}{Saber contemporáneo}
Excel calcula en un orden fijo. Primero lo que está entre \textbf{paréntesis}. Después las \textbf{potencias}. Después \textbf{multiplicar y dividir}, de izquierda a derecha. Al final \textbf{sumar y restar}, también de izquierda a derecha. Por eso \texttt{=2+3*4} da 14 y no 20: Excel multiplica antes de sumar. Si querías sumar primero, lo escribes: \texttt{=(2+3)*4}. La regla del oficio es corta. Cuando dudes del orden, pon paréntesis. No sobran nunca: le dicen a Excel, y a quien lea tu fórmula, qué se calcula primero. Una fórmula que otra persona entiende sin adivinar el orden es una fórmula bien escrita.
\end{softbox}

\begin{softbox}{milcMostaza}{milcGris}{Pregunta-puente}
\textit{La tejedora wounaan no puede cerrar la vuelta de afuera antes que la del centro. Cuando escribas \texttt{=B2*1+19~\%} para el IVA, ¿qué vuelta cerró Excel primero? ¿Era la que tú querías?}
\end{softbox}

\begin{softbox}{milcVerde}{milcGris}{Saber y saber hacer}
Al terminar podrás: (1) \textbf{identificar} en qué orden calcula Excel una fórmula con varios operadores; (2) \textbf{aplicar} paréntesis para que el orden sea el que necesitas; (3) \textbf{analizar} por qué una fórmula da un resultado distinto del esperado; (4) \textbf{evaluar} si una fórmula ajena se entiende sin adivinar el orden.
\end{softbox}



\small
\begin{tabularx}{\textwidth}{>{\bfseries}p{3.0cm}Y}
\rowcolor{milcGradePrimary!12}Autor & Dr. Álvaro Cárdenas Orozco\\
DBA & Tratamiento de datos cuantitativos con criterio ético (MEN, T\&I)\\
Referentes & Phronesis aristotélica · Floridi (ética del dato) · Estoicismo (ver lo que es)\\
Producto final & Una hoja de Excel con cinco fórmulas compuestas reales (nota final ponderada, precio con IVA y descuento, promedio sin el menor valor, entre otras), escritas con paréntesis que dejan claro el orden. Además, tres fórmulas mal escritas a propósito, corregidas y explicadas en una línea, y una prueba final con un cálculo de tu vida real.\\
\end{tabularx}
\normalsize

\puente{En la sesión 5 aprendiste a copiar fórmulas sin romperlas. Hoy aprendes a escribirlas cuando tienen varias operaciones. En la sesión 7 vas a graficar lo que calcules, y un gráfico sobre una fórmula mal ordenada engaña sin que nadie lo note.}

\milcestacion{milcGradeDark}{Ruta MILC: así vamos a aprender}
\begin{tabularx}{\textwidth}{>{\centering\arraybackslash}X>{\centering\arraybackslash}X>{\centering\arraybackslash}X>{\centering\arraybackslash}X}
\phasecard{1}{milcVerde}{milcGris}{Escucha\\[-1mm]\small Observo, escucho y pregunto.} &
\phasecard{2}{milcTurquesa}{milcGris}{Entender\\[-1mm]\small Sistematizo y ordeno.} &
\phasecard{3}{milcMagenta}{milcGris}{Hacer\\[-1mm]\small Construyo y aplico.} &
\phasecard{4}{milcVino}{milcGris}{Cierre\\[-1mm]\small Evalúo y reflexiono.}
\end{tabularx}


\puente{Antes de la teoría, vas a calcular tres expresiones a mano y después en Excel. Si algún resultado no coincide, ahí está la lección del día.}

\milcestacion{milcVerde}{Estación 1 de 3 · Escucha}

\begin{softbox}{milcVerde}{milcGris}{Escena para escuchar}
\textbf{Actividad 1 · IDENTIFICA --- A mano y en Excel} (15 min · individual). (1) Calcula a mano, con calculadora o en el cuaderno, estas tres expresiones: \texttt{2+3*4}, \texttt{(2+3)*4} y \texttt{20-6/2}. (2) Anota tus tres respuestas. (3) Abre Excel y escribe las tres como fórmulas, con el signo igual adelante. (4) Compara tus respuestas con las de Excel. (5) Si alguna no coincide, escribe cuál y qué crees que calculó Excel primero.
\end{softbox}

{\sffamily\bfseries\fontsize{8}{10}\selectfont\textcolor{milcNegro!55}{MARCA CUANDO LO HAYAS HECHO}}
\milccheckrow{Tengo las tres respuestas a mano y las tres de Excel.}
\milccheckrow{Sé cuál no coincidió, si alguna, y por qué.}
\milccheckrow{Puedo decir qué calcula Excel primero en 2+3*4.}

\begin{infoband}{milcVerde}


\textbf{Lo que va al cuaderno (Actividad 1 · IDENTIFICA).} \textbf{Título:} «Actividad 1 --- A mano y en Excel». \textbf{Formato:} tabla de 3 filas y 3 columnas (expresión / mi resultado / resultado de Excel). \textbf{Extensión:} un tercio de página. \textbf{Sabes que terminaste cuando} las tres filas están llenas y escribiste en una línea qué calcula Excel primero.
\end{infoband}

\puente{Ya viste que el orden cambia el resultado. Ahora entiendes los cuatro niveles del orden, y con tu pareja escribes cinco fórmulas reales con paréntesis que lo dejen claro.}

\milcestacion{milcTurquesa}{Estación 2 de 3 · Entender}

\begin{softbox}{milcTurquesa}{milcGris}{Lo que vas a entender}
\textbf{Actividad 2 · APLICA --- Cinco fórmulas reales con paréntesis} (30 min · parejas). (1) Con tu pareja, abran una hoja nueva y escriban tres notas parciales en B2, C2 y D2. (2) Escriban la nota final ponderada: \texttt{=(B2*0,3)+(C2*0,3)+(D2*0,4)}. Prueben con 3,0 en las tres: debe dar 3,0. (3) Escriban un precio en B5, un descuento en D5 y el IVA en E5. Calculen el precio final: \texttt{=B5*(1-D5)*(1+E5)}. (4) Escriban cinco valores en B8:F8 y calculen el promedio sin el menor: \texttt{=(SUMA(B8:F8)-MIN(B8:F8))/4}. (5) Inventen dos fórmulas más de su vida real, con al menos dos operadores y paréntesis, y pruébenlas con un caso que ya sepan.
\end{softbox}

\milcpaso{1}{milcTurquesa}{\textbf{Los cuatro niveles del orden.} Primero los paréntesis, de adentro hacia afuera. Segundo las potencias. Tercero multiplicar y dividir, de izquierda a derecha: \texttt{10/2*5} es 25, no 1. Cuarto sumar y restar, de izquierda a derecha.}
\milcpaso{2}{milcTurquesa}{\textbf{Los paréntesis dicen lo que pensaste.} \texttt{=B2*0,3+C2*0,3+D2*0,4} da bien por el orden de Excel. \texttt{=(B2*0,3)+(C2*0,3)+(D2*0,4)} da lo mismo y cualquiera lo lee. Cuando el orden que quieres no es el de Excel, los paréntesis son obligatorios. Cuando sí lo es, siguen ayudando a leer.}
\milcpaso{3}{milcTurquesa}{\textbf{Prueba con un caso que ya sabes.} Si la nota final con 3,0 en las tres notas no da 3,0, la fórmula está mal. Si el precio con 0~\% de descuento y 0~\% de IVA no da el precio base, la fórmula está mal. Un caso conocido delata el orden equivocado antes de que llegue a un gráfico.}
\milcpaso{4}{milcTurquesa}{\textbf{Léela en voz alta.} «Primero multiplico cada nota por su peso, después sumo las tres.» Si la frase sale limpia, la fórmula está bien ordenada. Si tienes que decir «bueno, en realidad Excel hace…», falta un paréntesis.}

\begin{drawbox}{milcTurquesa}{El orden de Excel, en una mirada}
\textbf{1. Paréntesis:} lo de adentro primero. \\[1mm] \textbf{2. Potencias.} \\[1mm] \textbf{3. Multiplicar y dividir:} de izquierda a derecha. \\[1mm] \textbf{4. Sumar y restar:} de izquierda a derecha. \\[1mm] \textbf{Duda:} pon paréntesis.
\end{drawbox}

\begin{softbox}{milcMostaza}{milcGris}{Errores comunes y su corrección}
(1) \textbf{Creer que se suma antes de multiplicar}: \texttt{=2+3*4} es 14. (2) \textbf{Restar en el numerador sin paréntesis}: \texttt{=A2-B2/C2} divide primero; para restar primero es \texttt{=(A2-B2)/C2}. (3) \textbf{IVA mal aplicado}: \texttt{=B2*1+19~\%} suma 0,19 al final; lo correcto es \texttt{=B2*(1+19~\%)}. (4) \textbf{Dividir entre lo que no es}: \texttt{=SUMA(B2:F2)-MIN(B2:F2)/4} divide solo el mínimo. (5) \textbf{No probar con un caso conocido}: el error llega intacto al gráfico.
\end{softbox}

\begin{infoband}{milcTurquesa}


\textbf{Lo que va al cuaderno (Actividad 2 · APLICA).} \textbf{Título:} «Actividad 2 --- Cinco fórmulas con paréntesis». \textbf{Formato:} tabla de 5 filas y 3 columnas (qué calcula / fórmula escrita / caso de prueba y resultado). \textbf{Extensión:} media página. \textbf{Sabes que terminaste cuando} las cinco fórmulas tienen su caso de prueba y las cinco dieron lo esperado.
\end{infoband}

\puente{Ya escribes fórmulas claras. Toca corregir las de otro: tres fórmulas mal escritas, una prueba con un cálculo de tu vida y la revisión de la fórmula de tu pareja.}

\milcestacion{milcMagenta}{Estación 3 de 3 · Hacer}

\begin{softbox}{milcMagenta}{milcGris}{Cómo vas a construir tu producto}
\textbf{Actividad 3 · EVALÚA --- Tres fórmulas rotas y una prueba real} (25 min · individual). (1) Copia estas tres fórmulas mal escritas: \texttt{=A2-B2/C2} cuando se quería restar primero; \texttt{=B2*1+19~\%} para el IVA; \texttt{=SUMA(B2:F2)-MIN(B2:F2)/4} para el promedio sin el menor. (2) Escribe al lado de cada una qué calculó Excel primero y por qué salió mal. (3) Corrígelas con paréntesis. (4) Haz una prueba real: un cálculo de tu vida con al menos dos operadores. Por ejemplo, el gasto del mes en transporte y tienda, o lo que queda de la mesada. (5) Lee la fórmula de tu pareja y di en voz alta qué se calcula primero; anota si acertaste.
\end{softbox}

\milcpaso{1}{milcMagenta}{\textbf{Tu entrega tiene tres piezas.} (1) Las cinco fórmulas de la Actividad 2 con su caso de prueba. (2) Las tres fórmulas corregidas, cada una con una línea que diga por qué fallaba. (3) La prueba real con su fórmula y su resultado.}
\milcpaso{2}{milcMagenta}{\textbf{Cómo se explica una fórmula rota.} «Excel dividió B2 entre C2 antes de restar, porque dividir va antes que restar. Yo quería restar primero. Solución: paréntesis en la resta.» Tres frases: qué hizo Excel, qué quería yo, qué cambié.}
\milcpaso{3}{milcMagenta}{\textbf{La prueba real vale más que el ejemplo.} Un cálculo de tu mesada, de tu transporte o de tu tiempo de pantalla te obliga a decidir el orden de verdad. Escribe primero la frase en palabras y después la fórmula.}
\milcpaso{4}{milcMagenta}{\textbf{Evaluar la fórmula de tu pareja.} Léela sin que te la explique y di qué se calcula primero. Si tuviste que adivinar, le falta un paréntesis. Díselo con la razón, no con el adjetivo.}

\begin{drawbox}{milcMagenta}{Antes de entregar}
\checkbox Cinco fórmulas con paréntesis y su caso de prueba. \\[1mm] \checkbox Tres fórmulas corregidas con una línea de explicación cada una. \\[1mm] \checkbox Una prueba real escrita primero en palabras y después en fórmula. \\[1mm] \checkbox Leí la fórmula de mi pareja y dije qué se calcula primero. \\[1mm] \checkbox Hoja guardada como G8-P1-S6-formulas.xlsx.
\end{drawbox}

\begin{softbox}{milcMagenta}{milcGris}{Plantilla de trabajo}
\textbf{Nota final.} \textit{=(B2*0,3)+(C2*0,3)+(D2*0,4)} \\[1mm] \textbf{Precio final.} \textit{=B5*(1-D5)*(1+E5)} \\[1mm] \textbf{Promedio sin el menor.} \textit{=(SUMA(B8:F8)-MIN(B8:F8))/4} \\[1mm] \textbf{Fórmula rota.} \textit{Qué hizo Excel, qué quería yo, qué cambié.} \\[1mm] \textbf{Prueba real.} \textit{Primero la frase, después la fórmula.}
\end{softbox}

\begin{infoband}{milcMagenta}


\textbf{Lo que va al cuaderno (Actividad 3 · EVALÚA).} \textbf{Título:} «Actividad 3 --- Tres fórmulas rotas y una prueba real». \textbf{Formato:} tabla de 3 filas y 3 columnas (fórmula rota / qué calculó Excel primero / fórmula corregida), más la prueba real en palabras y en fórmula. \textbf{Extensión:} media página. \textbf{Sabes que terminaste cuando} las tres corregidas dan el resultado esperado y tu pareja leyó tu prueba real sin adivinar el orden.
\end{infoband}

\puente{Lo que entregas hoy: cinco fórmulas con paréntesis, tres corregidas con su explicación y una prueba real. La prueba de calidad: tu pareja lee cada fórmula y dice en voz alta qué se calcula primero, sin equivocarse.}

\milcestacion{milcMostaza}{Producto de la sesión}
\textbf{Cinco fórmulas con paréntesis + tres corregidas y explicadas + una prueba real}.
\begin{softbox}{milcMostaza}{milcGris}{Criterios de calidad}
(1) Las cinco fórmulas tienen paréntesis que dejan claro el orden y un caso de prueba que dio bien. (2) Las tres fórmulas rotas están corregidas y cada una tiene una línea que dice por qué fallaba. (3) La prueba real está escrita primero en palabras y después en fórmula. (4) Tu pareja leyó una fórmula tuya y dijo qué se calcula primero sin adivinar. (5) La hoja está guardada con nombre claro.
\end{softbox}

\puente{Antes de cerrar, mira lo que hiciste desde cinco lados. Un paréntesis es una forma de respeto: le ahorra a quien lee la fórmula tener que adivinar lo que pensaste.}

\milcestacion{milcVino}{Evaluación liberadora · cinco dimensiones}

\begin{softbox}{milcVino}{milcGris}{Sentido de la evaluación}
Evaluar no es solo revisar una nota. Es reconocer qué aprendí, qué cambió en mí, qué emociones manejé, cómo me relaciono con la ciudad y la comunidad, y qué le dejaré a quien viene detrás.
\end{softbox}

\textbf{Mi reflexión liberadora en cinco dimensiones.} Completa cada frase iniciada:

\small
\begin{tabularx}{\textwidth}{>{\bfseries\RaggedRight\arraybackslash}p{3.4cm}Y>{\RaggedRight\arraybackslash}p{4.6cm}}
\rowcolor{milcVino}\color{white}Dimensión & \color{white}\bfseries Mi respuesta (frame starter) & \color{white}\bfseries Algo que descubrí\\
1. Personal & Lo que más cambió en mí fue \linead{6.5cm} & \linead{4.2cm}\\
2. Emocional & La emoción más fuerte fue \linead{2.5cm} y la regulé \linead{3cm} & \linead{4.2cm}\\
3. Ciudadana & En democracia esto importa porque \linead{6cm} & \linead{4.2cm}\\
4. Local (Cartago) & En mi barrio sería útil para \linead{6.5cm} & \linead{4.2cm}\\
5. Intergeneracional & A mi abuela/abuelo le diría \linead{6.5cm} & \linead{4.2cm}\\
\end{tabularx}
\normalsize

\begin{infoband}{milcVino}
\textbf{Para la próxima sustentación.} Vuelve a esta tabla en seis meses. Verás que las respuestas cambian — no porque lo aprendido cambie, sino porque tú cambias. La evaluación liberadora es una conversación contigo a través del tiempo.
\end{infoband}


\puente{Para terminar, tres ideas sobre hacer las cosas en orden, contadas con palabras de esta guía: Dussel sobre la tecnología como oficio que se piensa mientras se hace, Epicteto sobre decir solo lo necesario, Floridi sobre el orden por defecto que no siempre es el tuyo.}

\milcestacion{milcGradeDark}{Tres ideas para cerrar · Dussel, un estoico y Floridi}

\begin{softbox}{milcVino}{milcGris}{Enrique Dussel · lente del nosotros}
\textit{La tecnología no es teoría aplicada: es un oficio que se piensa mientras se hace.}

{\footnotesize\itshape\color{milcNegro!70}--- idea tomada de Enrique Dussel, contada con palabras de esta guía. No es una cita textual.}

Nadie aprende el orden de las operaciones de memoria y ya. Se aprende cuando una fórmula da 14 en vez de 20 y tienes que averiguar por qué. La tejedora tampoco aprende el rollo en un libro: lo aprende vuelta por vuelta.

\textbf{¿Qué entendí hoy porque una fórmula me salió mal, y no porque me lo dijeron?}
\end{softbox}
\linead{\linewidth}\\[1mm]
\linead{\linewidth}

\begin{softbox}{milcMostaza}{milcGris}{Estoicismo · Epicteto · cuidado interior}
\textit{Guarda silencio cuando puedas y di solo lo necesario, con pocas palabras.}

{\footnotesize\itshape\color{milcNegro!70}--- idea tomada de Epicteto, contada con palabras de esta guía. No es una cita textual.}

Un paréntesis es decir lo necesario: no adorna, aclara. Una fórmula con los paréntesis justos se lee de una vez. Una sin ellos obliga a adivinar. Una con paréntesis de más también cansa.

\textbf{¿Mis paréntesis dicen lo necesario, o puse de más porque no estaba seguro?}
\end{softbox}
\linead{\linewidth}\\[1mm]
\linead{\linewidth}

\begin{softbox}{milcTurquesa}{milcGris}{Luciano Floridi y la Onlife Initiative · lente de la infoesfera}
\textit{Las configuraciones por defecto de una tecnología deberían respetar a la persona que la usa.}

{\footnotesize\itshape\color{milcNegro!70}--- idea tomada de Luciano Floridi y la Onlife Initiative, contada con palabras de esta guía. No es una cita textual.}

El orden por defecto de Excel es una configuración: multiplica antes de sumar aunque tú no lo hayas pedido. Conocerlo es lo que te permite decidir tú, con paréntesis, en vez de dejar que la máquina decida por ti.

\textbf{¿En qué otra herramienta acepto un «por defecto» que no elegí?}
\end{softbox}
\linead{\linewidth}\\[1mm]
\linead{\linewidth}

\begin{infoband}{milcNegro}
\textbf{Nota para el docente.} Las tres ideas están contadas con palabras de esta guía a partir del pensamiento de cada autor; no son citas textuales. De 9.º en adelante se trabajan con la obra y su referencia.
\end{infoband}

\milcestacion{milcVino}{Mi autoevaluación · marca una casilla por fila}

{\small
\setlength{\extrarowheight}{2.2pt}
\begin{tabularx}{\textwidth}{>{\RaggedRight\arraybackslash\bfseries}p{3.0cm}YYY}
\rowcolor{milcVino}\color{white}\bfseries Criterio & \color{white}\bfseries Lo logré & \color{white}\bfseries En proceso & \color{white}\bfseries Necesito apoyo\\[.6mm]
Comprensión del tema & \milcopcion{Explico las ideas con mis palabras.} & \milcopcion{Reconozco las ideas, falta precisión.} & \milcopcion{Necesito repasar lo básico.}\\[.8mm]
\rowcolor{milcGris}Pensamiento computacional & \milcopcion{Usé los pilares al armar mi producto.} & \milcopcion{Usé algunos pilares.} & \milcopcion{Necesito releer la estación 2.}\\[.8mm]
Aplicación y producto & \milcopcion{Mi producto cumple los criterios.} & \milcopcion{Está armado, con ajustes pendientes.} & \milcopcion{Aún está en construcción.}\\[.8mm]
\rowcolor{milcGris}Reflexión 5 dimensiones & \milcopcion{Respondí cada una con honestidad.} & \milcopcion{Respondí algunas con detalle.} & \milcopcion{Mis respuestas son muy generales.}\\[.8mm]
Triángulo de pensamiento & \milcopcion{Conecté las 3 voces con mi trabajo.} & \milcopcion{Conecté al menos una.} & \milcopcion{No las relaciono con mi proceso.}\\[.8mm]
\end{tabularx}}

\vspace{3mm}
\textbf{Hoy entendí que:}\\[1mm]
\linead{\linewidth}\\[1mm]
\linead{\linewidth}

\textbf{Mi compromiso para la próxima sesión es:}\\[1mm]
\linead{\linewidth}\\[1mm]
\linead{\linewidth}

\begin{softbox}{milcMostaza}{milcGris}{Cita de despedida · Marco Aurelio}
\textit{``El obstáculo en el camino se vuelve el camino. Lo que estorba al camino se vuelve camino.''}\\[1mm]
Cada error de hoy, bien observado, se convierte en pieza del camino que pisarás mañana.
\end{softbox}

\small
\begin{tabularx}{\textwidth}{>{\bfseries}p{3.6cm}Y}
\rowcolor{milcGradeDark!12}Nombre completo: & \linead{10cm}\\
Fecha: & \linead{4cm} \quad Curso/grupo: \linead{4cm}\\
Mi firma: & \linead{9cm}\\
\end{tabularx}
\normalsize

\begin{infoband}{milcGradeDark}
\textbf{Para la próxima sesión.} Llega con tus cinco fórmulas, las tres corregidas y la prueba real. En la sesión 7 vas a hacer gráficos de barras, de líneas y circulares, y a aprender cuándo un gráfico engaña. Lo que calculaste hoy es lo que vas a graficar.
\end{infoband}

\end{document}
